\section{CRT alias compression and its non-tensor boundary}
\label{sec:tpc35-crt-alias}

The single--modulus conflict in TPC-33 is elementary: orbit coverage
prefers \(q\lesssim J\), whereas no--wrap on an affine determinant of
size \(X\) requires \(q\gg X\).  Several small moduli can increase the
joint modulus without sacrificing local orbit coverage.  We now compute
the exact threshold and then prove why independent one--coordinate bounds
do not automatically multiply.

\subsection{Exact alias identity}

\begin{theorem}[Finite CRT alias decomposition]
\label{thm:tpc35-alias-decomposition}
Let \(q_1,\ldots,q_r\) be positive integers, let
\begin{equation}
 M=[q_1,\ldots,q_r],
 \label{eq:tpc35-joint-modulus}
\end{equation}
and let \(F\in\Z\) satisfy \(|F|\le B\).  Then
\begin{equation}
 \boxed{
 \prod_{i=1}^r\one_{q_i\mid F}
 =\one_{M\mid F}
 =\sum_{|k|\le\lfloor B/M\rfloor}\one_{F=kM}.}
 \label{eq:tpc35-alias-identity}
\end{equation}
For all integers \(F\) in the ambient interval \(|F|\le B\), the joint
congruence is equivalent to \(F=0\) if and only if \(M>B\).
\end{theorem}

\begin{proof}
Divisibility by every \(q_i\) is equivalent to divisibility by their least
common multiple \(M\).  Every multiple \(F=kM\) in the stated interval
has \(|k|\le\lfloor B/M\rfloor\), proving the identity.  If \(M>B\), the
only possible \(k\) is zero.  Conversely, if \(M\le B\), the value \(F=M\)
is a nonzero counterexample to a uniform no--wrap assertion on the whole
ambient interval.
\end{proof}

Applied to
\begin{equation}
 F=sU-aD-h,
 \label{eq:tpc35-affine-determinant}
\end{equation}
whose physical ambient range is \(B\asymp X\), this is an exact
decomposition of a congruence relaxation into parallel integer determinant
fibers.  It does not bound the mass on any one fiber.

\begin{corollary}[The orbit-scale alias ladder]
\label{cor:tpc35-alias-ladder}
Suppose \(q_1,\ldots,q_r\) are pairwise coprime and
\(q_i\asymp J=X^{133/400+o(1)}\).  Then the number of possible \(k\) in
\eqref{eq:tpc35-alias-identity} is
\begin{equation}
 O\!\left(1+\frac{X}{J^r}\right),
 \qquad
 \frac{X}{J^r}=X^{1-r(133/400)+o(1)}.
 \label{eq:tpc35-alias-count}
\end{equation}
In particular,
\begin{center}
\begin{tabular}{@{}ccc@{}}
\toprule
\(r\) & \(X/J^r\) & Interpretation \\
\midrule
1 & \(X^{267/400+o(1)}\) & one-modulus alias family \\
2 & \(X^{134/400+o(1)}\) & two-modulus alias family \\
3 & \(X^{1/400+o(1)}\) & first diagonal-budget scale \\
4 & \(X^{-132/400+o(1)}\) & strict no-wrap for large \(X\) \\
\bottomrule
\end{tabular}
\end{center}
Among orbit-scale choices \(q_i\asymp J\), four moduli are the first
that force strict no--wrap on an \(O(X)\) ambient determinant interval.
\end{corollary}

Because \(QJ=X^{1+o(1)}\), the three--modulus endpoint is
\begin{equation}
 \boxed{
 \frac{X}{J^3}=\frac{Q}{J^2}=X^{1/400+o(1)}.}
 \label{eq:tpc35-critical-three-modulus}
\end{equation}
This is the same factor by which the allowed orbit energy \(Q^3/J\)
exceeds the proved diagonal \(Q^2J\).  Hence the following accounting
statement is exact: if a joint three--modulus decomposition has
\(O(X^{o(1)}X/J^3)\) alias fibers and every fiber can be reassembled at
scale \(O(X^{o(1)}Q^2J)\), then absolute reassembly gives
\begin{equation}
 \frac{X}{J^3}\,Q^2J
 =\frac{Q^3}{J}.
 \label{eq:tpc35-conditional-alias-reassembly}
\end{equation}
Neither the fiber estimate nor its uniformity is proved here.  In
particular, \eqref{eq:tpc35-critical-three-modulus} is exponent
compatibility, not closure of the physical gate.

\subsection{The exact two-congruence compatibility condition}

We record a small point needed to avoid a false CRT simplification.
TPC-19 contains the primitive determinant--divisor normal form, while
TPC-29 proves the general compatibility \(g\mid h(m_1-m_2)\), the unique
lcm class, and its primitive specialization
\citep{WangTPC19,WangTPC29}.  The self-contained restatement below makes
the reduced divisor \(g/(g,h)\) explicit and records a useful
counterexample; the compatibility principle itself is not new.

\begin{proposition}[Affine CRT compatibility]
\label{prop:tpc35-affine-CRT-compatibility}
Let \(u_1,u_2\ge1\), let \(g=(u_1,u_2)\), and assume
\((m_i,u_i)=1\).  Then the congruences
\begin{equation}
 m_i j\equiv-h\pmod{u_i},\qquad i=1,2,
 \label{eq:tpc35-two-congruences}
\end{equation}
are compatible if and only if
\begin{equation}
 \boxed{g\mid h(m_1-m_2),}
 \label{eq:tpc35-general-compatibility}
\end{equation}
equivalently
\begin{equation}
 \frac{g}{(g,h)}\mid m_1-m_2.
 \label{eq:tpc35-reduced-compatibility}
\end{equation}
When compatible, they determine one residue class modulo
\([u_1,u_2]\).  Under the physical primitive hypothesis
\begin{equation}
 (u_i,hm_i)=1,\qquad i=1,2,
 \label{eq:tpc35-physical-primitive-CRT}
\end{equation}
compatibility reduces to
\begin{equation}
 \boxed{(u_1,u_2)\mid m_1-m_2.}
 \label{eq:tpc35-primitive-compatibility}
\end{equation}
\end{proposition}

\begin{proof}
The individual solutions are
\[
 j\equiv-hm_i^{-1}\pmod{u_i}.
\]
The generalized CRT says they are compatible exactly when the two
residues agree modulo \(g\):
\[
 hm_1^{-1}\equiv hm_2^{-1}\pmod g.
\]
Multiplication by \(m_1m_2\), which is invertible modulo \(g\), gives
\eqref{eq:tpc35-general-compatibility}.  Dividing by the common factor
\((g,h)\) gives \eqref{eq:tpc35-reduced-compatibility}.  Under
\eqref{eq:tpc35-physical-primitive-CRT}, \((g,h)=1\), giving
\eqref{eq:tpc35-primitive-compatibility}.  The modulus of the joint class
is the least common multiple by the generalized CRT.
\end{proof}

The condition \((m_i,hu_i)=1\) is not a substitute for
\eqref{eq:tpc35-physical-primitive-CRT}: it makes \(m_i\) invertible but
need not make \(h\) invertible modulo \(u_i\).  For example,
\(h=2,u_1=u_2=4,m_1=1,m_2=3\) gives compatible congruences although
\(4\nmid m_1-m_2\); the general criterion
\eqref{eq:tpc35-general-compatibility} is correct.

\subsection{CRT lifts do not force tensor gain}

Let \(M=M_0M'\), where \((M_0,M')=1\).  For
\(f,g:\Z/M_0\Z\to\C\), lift them to \(\Z/M\Z\) by
\begin{equation}
 F(D)=f(D\bmod M_0),\qquad G(U)=g(U\bmod M_0).
 \label{eq:tpc35-coordinate-lifts}
\end{equation}
For units \(a,s\pmod M\), put
\begin{equation}
 \cC_M(a,s)
 :=\sum_{D\bmod M}F(D)
 G\!\left(s^{-1}(aD+h)\right).
 \label{eq:tpc35-ring-correlation}
\end{equation}
Let \(\bar F=M^{-1}\sum_DF(D)\) and similarly for all means.

\begin{theorem}[Proper-coordinate non-tensorization]
\label{thm:tpc35-CRT-lift}
With \(a_0,s_0,h_0\) denoting reduction modulo \(M_0\), one has
\begin{equation}
 \boxed{
 \cC_M(a,s)-M\bar F\bar G
 =M'\bigl(\cC_{M_0}(a_0,s_0)-M_0\bar f\bar g\bigr).}
 \label{eq:tpc35-CRT-pointwise-lift}
\end{equation}
Consequently
\begin{align}
 &\boxed{
 \sum_{a,s\in(\Z/M\Z)^\times}
 |\cC_M(a,s)-M\bar F\bar G|^2}
 \notag\\[-0.2em]
 &\boxed{\qquad
 =M'^2\varphi(M')^2
 \sum_{a_0,s_0\in(\Z/M_0\Z)^\times}
 |\cC_{M_0}(a_0,s_0)-M_0\bar f\bar g|^2.}
 \label{eq:tpc35-CRT-energy-lift}
\end{align}
In the normalized rms
\begin{equation}
 \operatorname{rms}_M^2
 :=\frac1{\varphi(M)^2}
 \sum_{a,s\in(\Z/M\Z)^\times}
 \left|\frac{\cC_M(a,s)-M\bar F\bar G}{M}\right|^2,
 \label{eq:tpc35-normalized-ring-rms}
\end{equation}
the value is exactly \(\operatorname{rms}_{M_0}^2\).
\end{theorem}

\begin{proof}
Every residue \(D_0\pmod {M_0}\) has exactly \(M'\) lifts modulo \(M\), while
the lifted summand depends only on \(D_0,a_0,s_0\).  Hence
\[
 \cC_M(a,s)=M'\cC_{M_0}(a_0,s_0),
 \qquad \bar F=\bar f,\quad\bar G=\bar g,
\]
which proves \eqref{eq:tpc35-CRT-pointwise-lift}.  Each pair
\((a_0,s_0)\) has \(\varphi(M')^2\) unit lifts.  Squaring the pointwise
identity and summing gives \eqref{eq:tpc35-CRT-energy-lift}.  Finally use
\(
 \varphi(M)=\varphi(M_0)\varphi(M')
\)
and \(M=M_0M'\) in \eqref{eq:tpc35-normalized-ring-rms}; all \(M'\)-factors
cancel.
\end{proof}

Taking \(M_0\) to be the product of any proper subset of the prime
coordinates shows that every proper-coordinate face of a squarefree CRT
product may retain the normalized rms of the smaller product.  Extra
coordinates only repeat data that do not depend on them.
This is the orbit--scale form of the low-coordinate phenomenon already
encountered in earlier conductor decompositions
\citep{WangTPC22,WangTPC24}.  A joint three--modulus proof must control
principal, singleton, and double-coordinate modes; the full-coordinate
piece alone does not suffice.

\subsection{Alias--coverage bookkeeping for a merged modulus}

One might merge small moduli into \(P\), use
\(O(1+X/P)\) aliases, and then complete the \(J\)-long orbit modulo \(P\),
at cost \(O(1+P/J)\).  This specified route obeys an exact bookkeeping
inequality.

\begin{proposition}[Merge--then--complete cost ledger]
\label{prop:tpc35-alias-coverage-law}
For every \(P,X,J>0\),
\begin{equation}
 \boxed{
 \left(1+\frac XP\right)
 \left(1+\frac PJ\right)
 \ge\frac XJ.}
 \label{eq:tpc35-alias-coverage-law}
\end{equation}
When \(X\asymp QJ\), the right side is \(\asymp Q\).
\end{proposition}

\begin{proof}
Expanding the left side produces the nonnegative terms
\(1+X/P+P/J\) in addition to \(X/J\).
\end{proof}

At \(P\asymp J^3\), the alias factor reaches the critical
\(X/J^3=X^{1/400+o(1)}\), but the completion factor is \(J^2\); their
product still contains \(X/J\asymp Q\).  This proposition is not a no--go
theorem for a joint delta method, circle method, multi--modulus
large--sieve estimate, or tensor spectral estimate.  It rules out only
the strategy that
first forgets the separate CRT coordinates and then pays for complete
coverage of the merged modulus.
